\chapter{Thực nghiệm, kết quả và đánh giá}
\label{chap:thuc_nghiem}

Chương này trình bày toàn diện các kết quả thực nghiệm, đối sánh định lượng và kiểm định thống kê chuẩn mực cho cả hai giai đoạn của hệ thống Two-Stage Pipeline trên bộ dữ liệu thực tế IC4SD-Traffic-HCM: Giai đoạn 1 (ước lượng số lượng phương tiện từ ảnh thị giác) và Giai đoạn 2 (dự báo lưu lượng không -- thời gian trên mạng lưới đồ thị 608 nút camera). Trọng tâm phân tích được đặt vào cơ chế \textbf{lan truyền sai số nhận dạng (Error Propagation)} và \textbf{năng lực kháng nhiễu cảm biến (Noise Robustness)} giữa các dòng kiến trúc học sâu không -- thời gian.

% ==============================================================================
\section{Thiết lập môi trường và tiêu chí đánh giá thực nghiệm}
\label{sec:thiet_lap_moi_truong}

\subsection{Môi trường phần cứng và phần mềm}
Các mô hình trong nghiên cứu được cài đặt trên nền tảng PyTorch 2.x và PyTorch Geometric. Quá trình huấn luyện và đánh giá được thực thi trên hệ thống máy chủ tính toán hiệu năng cao trang bị \textbf{2 GPU NVIDIA RTX 4090 (24GB VRAM mỗi GPU)}, bộ vi xử lý AMD EPYC 7543 32 lõi và 128GB RAM DDR4. Quá trình huấn luyện mô hình đồ thị áp dụng cơ chế song song hóa dữ liệu (DataParallel) phân bổ đều 32 mẫu trên mỗi GPU (tổng batch size $B=64$).

\subsection{Đặc tả cấu trúc tham số và bộ nhớ kích hoạt GPU của các mô hình Stage 2}
Bảng \ref{tab:sieu_tham_so_chap5_models} đối sánh cấu trúc tham số, cơ chế xử lý không gian -- thời gian và mức tiêu thụ bộ nhớ kích hoạt GPU đỉnh (Peak GPU Activation Memory) của 8 mô hình dự báo không -- thời gian được đánh giá trong Giai đoạn 2.

\begin{table}[htbp]
\centering
\caption{Đối sánh kiến trúc, tham số và bộ nhớ GPU đỉnh trong Giai đoạn 2 ($N=608$ nút, $T_{\text{in}}=24, H=6$).}
\label{tab:sieu_tham_so_chap5_models}
\small
\resizebox{\textwidth}{!}{%
\begin{tabular}{lcccccc}
\toprule
\textbf{Kiến trúc mô hình} & \textbf{Số tham số} & \textbf{RAM đỉnh (MB)} & \textbf{Toán tử Không gian} & \textbf{Toán tử Thời gian} & \textbf{Số đầu ($h$)} & \textbf{Số khối} \\
\midrule
STGCN Baseline \cite{stgcn} & 375.212 & 4.837,6 & ChebNet ($K=3$) & 1D Gated GLU & 0 & 2 Khối \\
STGCN (Large) & 487.420 & 5.793,9 & ChebNet ($K=3$) & 1D Gated GLU & 0 & 2 Khối \\
Graph WaveNet \cite{wavenet} & 548.620 & 7.731,4 & Adaptive GCN & Dilated Causal Conv & 0 & 4 Khối \\
ASTGCN \cite{astgcn} & 1.532.566 & 9.260,7 & ChebNet + Spatial Attn & 1D Conv + Temporal Attn & Không/Thời gian & 2 Khối \\
STAEformer \cite{staeformer} & 568.492 & 11.509,1 & Adaptive Spatial Emb & Spatio-Temporal Trans & $h=4$ & 1 Khối \\
MegaCRN \cite{megacrn} & 527.402 & 4.850,4 & Meta-Graph Emb & GRU Enc-Dec & 0 & GRU Layers \\
DSTAGNN \cite{dstagnn} & 910.796 & 9.801,5 & Dyn. Softmax Graph & Temporal Transformer & $h=4$ & 1 Khối \\
iTransformer \cite{itransformer} & \textbf{149.516} & \textbf{265,2} & Variate Attention & Inverted MLP & $h=8$ & 2 Tầng \\
\midrule
\textbf{TA-STGCN (Đề xuất)} & \textbf{453.052} & \textbf{8.606,9} & ChebNet ($K=3$) & 1D GLU + Temporal Self-Attn & $h=4$ & 2 Khối \\
\bottomrule
\end{tabular}%
}
\end{table}

\subsection{Cấu hình siêu tham số huấn luyện}
Bảng \ref{tab:sieu_tham_so_chap5} tổng hợp cấu hình siêu tham số huấn luyện chuẩn mực được áp dụng đồng nhất cho Giai đoạn 1 và Giai đoạn 2.

\begin{table}[htbp]
\centering
\caption{Cấu hình siêu tham số huấn luyện chi tiết cho Giai đoạn 1 và Giai đoạn 2.}
\label{tab:sieu_tham_so_chap5}
\small
\begin{tabularx}{\textwidth}{l X X}
\toprule
\textbf{Siêu tham số} & \textbf{Giai đoạn 1 (Nhận dạng thị giác)} & \textbf{Giai đoạn 2 (Dự báo đồ thị)} \\
\midrule
Thuật toán tối ưu hóa & AdamW & Adam \\
Tốc độ học khởi tạo ($\eta$) & $3 \times 10^{-4}$ & $1 \times 10^{-3}$ \\
Hệ số suy giảm trọng số (Weight Decay) & $1 \times 10^{-2}$ & $1 \times 10^{-4}$ \\
Chiến lược điều chỉnh tốc độ học & Cosine Annealing ($\eta_{\min} = 10^{-6}$) & ReduceLROnPlateau ($\gamma=0.5$, patience=10) \\
Số epoch tối đa & 120 epochs & 120 epochs \\
Ngưỡng dừng sớm (Early Stopping) & 20 epochs & 20 epochs \\
Kích thước mini-batch & 32 & 64 (32 mẫu / GPU) \\
Hàm mất mát & Smooth $L_1$ (Huber, $\delta=1.0$) & Huber Loss ($\delta=1.0$, dùng chung) \\
Tăng cường dữ liệu & Lật ngang, Color Jitter, Random Erase & Không (Bảo toàn trật tự chuỗi thời gian) \\
\bottomrule
\end{tabularx}
\end{table}

\subsection{Tiêu chí đánh giá và quy chuẩn báo cáo}
Để phản ánh trung thực năng lực vận hành thực tế, toàn bộ các chỉ số sai số được tính toán trên \textbf{giá trị số lượng phương tiện thực tế sau khi đã giải chuẩn hóa} thông qua biến đổi nghịch đảo: $y_{\text{real}} = y_{\text{z}} \cdot \sigma_i + \mu_i$, trong đó $\mu_i$ và $\sigma_i$ là kỳ vọng và độ lệch chuẩn của từng nút tính toán cố định trên tập huấn luyện (Train).

Ba thước đo sai số chuẩn mực được sử dụng gồm:
\begin{align}
    \text{MAE} &= \frac{1}{N \cdot H \cdot M} \sum_{i=1}^{N} \sum_{h=1}^{H} \sum_{k=1}^{M} \left| y_{i,h}^{(k)} - \hat{y}_{i,h}^{(k)} \right| \\
    \text{RMSE} &= \sqrt{\frac{1}{N \cdot H \cdot M} \sum_{i=1}^{N} \sum_{h=1}^{H} \sum_{k=1}^{M} \left( y_{i,h}^{(k)} - \hat{y}_{i,h}^{(k)} \right)^2} \\
    \text{MAPE} &= \frac{1}{|\Omega|} \sum_{(i,h,k) \in \Omega} \frac{\left| y_{i,h}^{(k)} - \hat{y}_{i,h}^{(k)} \right|}{y_{i,h}^{(k)} + \epsilon_{\mathrm{MAPE}}} \times 100\%
\end{align}
trong đó $N = 608$ nút mạng, $H = 6$ bước thời gian dự báo, $M$ là số mẫu chuỗi kiểm thử, $\Omega = \{(i,h,k) \mid y_{i,h}^{(k)} > 0.5\}$ là tập các điểm quan trắc có dòng xe thực tế lưu thông nhằm tránh sai số vô nghĩa khi lưu lượng bằng 0, và $\epsilon_{\mathrm{MAPE}} = 10^{-5}$ là hằng số chống chia cho 0.

Toàn bộ các thí nghiệm được lặp lại độc lập trên \textbf{5 seed ngẫu nhiên} $\mathcal{S} = \{42, 100, 2024, 22, 99\}$. Kết quả báo cáo theo định dạng thống kê chuẩn mực $\text{Mean} \pm \text{Std}$ cùng với khoảng tin cậy Bootstrap 95\% ($B = 1.000$ lần lấy mẫu lại có hoàn lại).

% ==============================================================================
\section{Kết quả thực nghiệm Giai đoạn 1 (Vision Perception Benchmark)}
\label{sec:thuc_nghiem_giai_doan_1}

\subsection{Đối sánh định lượng các kiến trúc mạng xương sống}
Bảng \ref{tab:stage1_benchmark_full_chap5} trình bày kết quả đánh giá toàn diện 5 kiến trúc thị giác máy tính tiêu biểu trên 501 ảnh kiểm thử độc lập từ camera giao thông TP. Hồ Chí Minh qua 5 seed ngẫu nhiên.

\begin{table}[htbp]
\centering
\caption{Đánh giá toàn diện các mô hình đếm xe Stage 1 qua 5 seed ngẫu nhiên ($\text{Mean} \pm \text{Std}$).}
\label{tab:stage1_benchmark_full_chap5}
\small
\resizebox{\textwidth}{!}{%
\begin{tabular}{lccccc}
\toprule
\textbf{Chỉ số đánh giá} & \textbf{EfficientNet-B5} & \textbf{ResNet-50} & \textbf{MobileNetV3-Large} & \textbf{ViT-Base/16} & \textbf{ConvNeXt-Tiny} \\
\midrule
Số tham số (Params) & 28.3M & 23.5M & \textbf{4.2M} & 85.8M & 27.8M \\
Độ phức tạp FLOPs & 75.3 GFLOPs & 132.2 GFLOPs & \textbf{6.9 GFLOPs} & 539.1 GFLOPs & 142.6 GFLOPs \\
Độ trễ suy luận (ms/batch 32) & $35.2 \pm 1.0$ & $34.6 \pm 0.7$ & $\mathbf{34.3 \pm 0.5}$ & $45.0 \pm 0.8$ & $34.9 \pm 0.9$ \\
\midrule
\textbf{MAE Tổng thể} & $\mathbf{3.6363 \pm 0.0619}$ & $3.6840 \pm 0.0534$ & $4.0126 \pm 0.0578$ & $5.5783 \pm 0.3120$ & $6.2903 \pm 3.6079$ \\
95\% Bootstrap CI MAE & $[3.5595 - 3.7131]$ & $[3.6177 - 3.7503]$ & $[3.9408 - 4.0844]$ & $[5.1908 - 5.9657]$ & $[1.8105 - 10.7701]$ \\
MAPE Tổng thể (\%) & $\mathbf{43.84\% \pm 1.20\%}$ & $45.83\% \pm 0.26\%$ & $52.05\% \pm 2.21\%$ & $68.79\% \pm 6.09\%$ & $80.11\% \pm 48.03\%$ \\
RMSE Tổng thể & $\mathbf{5.5345 \pm 0.0781}$ & $5.5745 \pm 0.0705$ & $6.0903 \pm 0.0756$ & $8.3289 \pm 0.5077$ & $9.6716 \pm 5.6654$ \\
\midrule
MAE Xe ô tô con & $\mathbf{2.1723 \pm 0.0415}$ & $2.2240 \pm 0.0642$ & $2.4679 \pm 0.0433$ & $3.3948 \pm 0.2252$ & $3.3374 \pm 1.5211$ \\
MAE Xe máy & $\mathbf{5.1003 \pm 0.0969}$ & $5.1440 \pm 0.0448$ & $5.5572 \pm 0.1267$ & $7.7617 \pm 0.4021$ & $9.2432 \pm 5.6963$ \\
\midrule
Hệ số biến thiên phương sai (CV) & $1.70\%$ & $1.45\%$ & $\mathbf{1.44\%}$ & $5.59\%$ & $57.36\%$ \\
\bottomrule
\end{tabular}%
}
\end{table}

\begin{figure}[htbp]
\centering
\includegraphics[width=0.88\linewidth]{fig/stage1_density_mae.pdf}
\caption{So sánh sai số MAE của các kiến trúc thị giác Stage 1 theo từng phân tầng mật độ phương tiện.}
\label{fig:stage1_density_mae_chap5}
\end{figure}

\subsection{Phân tích phương sai khởi tạo và độ ổn định tối ưu hóa}
Bảng \ref{tab:stage1_seed_variance} chi tiết hóa phương sai khởi tạo ($\sigma^2$), độ lệch chuẩn ($\sigma$), khoảng tin cậy Bootstrap 95\% và hệ số biến thiên ($\text{CV} = \sigma / \mu \times 100\%$) của chỉ số MAE tổng thể trên 5 seed ngẫu nhiên.

\begin{table}[htbp]
\centering
\caption{Thống kê phương sai và độ ổn định khởi tạo của các mô hình Stage 1 qua 5 seed ngẫu nhiên.}
\label{tab:stage1_seed_variance}
\small
\begin{tabular}{lcccc}
\toprule
\textbf{Kiến trúc Backbone} & \textbf{Mean MAE} & \textbf{95\% Bootstrap CI} & \textbf{Phương sai ($\sigma^2$)} & \textbf{CV (\%)} \\
\midrule
MobileNetV3-Large \cite{mobilenet} & $4.0126$ & $3.9408 - 4.0844$ & $3.34 \times 10^{-3}$ & $\mathbf{1.44\%}$ \\
ResNet-50 \cite{resnet} & $3.6840$ & $3.6177 - 3.7503$ & $2.85 \times 10^{-3}$ & $1.45\%$ \\
EfficientNet-B5 \cite{efficientnet} & $\mathbf{3.6363}$ & $\mathbf{3.5595 - 3.7131}$ & $3.83 \times 10^{-3}$ & $1.70\%$ \\
ViT-Base/16 \cite{vit} & $5.5783$ & $5.1908 - 5.9657$ & $9.74 \times 10^{-2}$ & $5.59\%$ \\
ConvNeXt-Tiny \cite{convnext} & $6.2903$ & $1.8105 - 10.7701$ & $1.30 \times 10^{1}$ & $57.36\%$ \\
\bottomrule
\end{tabular}
\end{table}

\textbf{Luận giải khoa học từ kết quả Giai đoạn 1:}
\begin{enumerate}
    \item \textbf{Hiệu năng vượt trội và tính kinh tế của EfficientNet-B5:} EfficientNet-B5 đạt độ chính xác cao nhất trên mọi tiêu chí (MAE tổng thể đạt $3.6363 \pm 0.0619$, MAE ô tô đạt $2.1723$, MAE xe máy đạt $5.1003$). Đáng chú ý, EfficientNet-B5 chỉ tiêu tốn 75.3 GFLOPs, tiết kiệm $43.0\%$ lượng phép tính so với ResNet-50 (132.2 GFLOPs) trong khi đạt độ trễ suy luận tương đương ($35.2$ ms so với $34.6$ ms trên batch 32). Kết quả này khẳng định tính ưu việt của nguyên lý mở rộng đồng thời (Compound Scaling), xác lập EfficientNet-B5 là mạng xương sống phù hợp nhất để sinh nhãn dữ liệu cho Giai đoạn 2.
    \item \textbf{Giải pháp siêu nhẹ cho phần cứng biên (MobileNetV3-Large):} MobileNetV3-Large thể hiện tính kinh tế ấn tượng: chỉ với 4.2M tham số và 6.9 GFLOPs (nhẹ hơn 6.7 lần về tham số và gần 11 lần về FLOPs so với EfficientNet-B5), mô hình vẫn đạt sai số chấp nhận được là $\text{MAE} = 4.0126 \pm 0.0578$ với độ ổn định khởi tạo cao nhất ($\text{CV} = 1.44\%$). Đây là cấu hình thay thế khả thi cho các cụm camera biên có kinh phí phần cứng hạn hẹp.
    \item \textbf{Độ nhạy cảm của các kiến trúc thuần Attention khi tinh chỉnh trên dữ liệu vừa:} Các kiến trúc CNN truyền thống tích hợp Batch Normalization (MobileNetV3, ResNet-50, EfficientNet-B5) duy trì tính hội tụ cực kỳ ổn định với phương sai rất nhỏ ($\sigma^2 < 3.84 \times 10^{-3}$, $\text{CV} \le 1.70\%$). Ngược lại, ViT-Base/16 có phương sai lớn hơn ($\text{CV} = 5.59\%$), trong khi ConvNeXt-Tiny thể hiện sự dao động mạnh giữa các lần chạy ngẫu nhiên ($\text{CV} = 57.36\%, \sigma^2 = 13.017$). Điều này phản ánh rằng các kiến trúc sử dụng Layer Normalization và kernel tích chập lớn $7 \times 7$ nhạy cảm hơn với trọng số khởi tạo khi thích nghi vào bài toán hồi quy liên tục trên tập dữ liệu kích thước vừa.
    \item \textbf{Kiểm định ý nghĩa thống kê:} Kiểm định phi tham số Friedman omnibus cho thấy sự phân hóa hiệu năng giữa 5 mạng xương sống có ý nghĩa thống kê cực kỳ rõ rệt ($\chi^2 = 863.55, p < 0.001$). Kiểm định Wilcoxon Signed-Rank hậu nghiệm kèm hiệu chỉnh Holm-Bonferroni xác nhận rằng EfficientNet-B5 vượt trội có ý nghĩa thống kê so với ViT-Base/16 ($p < 0.001$, Cohen's $d = 0.522$), ConvNeXt-Tiny ($p < 0.001$, Cohen's $d = 0.447$) và MobileNetV3-Large ($p < 0.001$, Cohen's $d = 0.164$). Khoảng cách sai số giữa EfficientNet-B5 và ResNet-50 không có ý nghĩa thống kê ($p = 0.253$, Cohen's $d = -0.022$), khẳng định cả hai cấu trúc CNN đều đạt độ chính xác hàng đầu trong bài toán nhận dạng.
\end{enumerate}

% ==============================================================================
\section{Phân tích sai số nhận dạng theo mật độ giao thông và kiểm chứng Grad-CAM++}
\label{sec:phan_tich_sai_so_theo_mat_do}

\subsection{Độ chệch nhận dạng phân tầng theo mật độ}
Để đánh giá tác động của mật độ phương tiện tới sai số thị giác, tập kiểm thử gồm 501 ảnh được phân tầng thành ba mức độ hoạt động: Mật độ Thấp ($<10$ xe/ảnh, 103 ảnh), Mật độ Trung bình ($10-25$ xe/ảnh, 151 ảnh) và Mật độ Cao ($>25$ xe/ảnh, 247 ảnh). Bảng \ref{tab:stage1_density_bias_chap5} chi tiết hóa độ chệch dự đoán ($\text{Bias} = \hat{y} - y$) theo phân lớp phương tiện trên từng phân tầng.

\begin{table}[htbp]
\centering
\caption{Độ chệch nhận dạng Stage 1 phân tầng theo mật độ giao thông ($\text{Mean} \pm \text{Std}$).}
\label{tab:stage1_density_bias_chap5}
\small
\resizebox{\textwidth}{!}{%
\begin{tabular}{llccc}
\toprule
\textbf{Kiến trúc Backbone} & \textbf{Phân loại} & \textbf{Mật độ Thấp ($<10$)} & \textbf{Mật độ TB ($10-25$)} & \textbf{Mật độ Cao ($>25$)} \\
\midrule
\multirow{2}{*}{\textbf{EfficientNet-B5}} & Ô tô con & $-0.03 \pm 0.24$ & $-0.30 \pm 0.30$ & $\mathbf{-0.24 \pm 0.50}$ \\
 & Xe máy & $+0.51 \pm 0.74$ & $+0.31 \pm 0.56$ & $\mathbf{-0.63 \pm 0.87}$ \\
\midrule
\multirow{2}{*}{ResNet-50} & Ô tô con & $+0.14 \pm 0.17$ & $-0.05 \pm 0.11$ & $-0.13 \pm 0.42$ \\
 & Xe máy & $+0.43 \pm 0.53$ & $+0.12 \pm 0.54$ & $-0.70 \pm 1.20$ \\
\midrule
\multirow{2}{*}{MobileNetV3-Large} & Ô tô con & $-0.03 \pm 0.09$ & $-0.33 \pm 0.23$ & $-0.37 \pm 0.74$ \\
 & Xe máy & $+0.56 \pm 0.31$ & $+0.14 \pm 0.56$ & $-0.85 \pm 1.66$ \\
\midrule
\multirow{2}{*}{ViT-Base/16} & Ô tô con & $+0.51 \pm 0.32$ & $-0.30 \pm 0.25$ & $-1.20 \pm 0.48$ \\
 & Xe máy & $+1.97 \pm 1.00$ & $+0.99 \pm 0.93$ & $-4.63 \pm 0.80$ \\
\midrule
\multirow{2}{*}{ConvNeXt-Tiny} & Ô tô con & $+1.31 \pm 1.54$ & $-0.26 \pm 0.32$ & $-2.63 \pm 3.22$ \\
 & Xe máy & $+5.57 \pm 6.91$ & $+1.80 \pm 2.36$ & $-10.22 \pm 13.11$ \\
\bottomrule
\end{tabular}%
}
\end{table}

\textbf{Phân tích hành vi độ chệch:}
\begin{itemize}
    \item Tại mật độ thấp ($<10$ xe), tất cả các mô hình xuất hiện độ chệch dương nhẹ đối với xe máy ($+0.43$ đến $+0.56$), do mạng thỉnh thoảng phản ứng nhầm với các chi tiết biển quảng cáo hoặc vệt sơn trên mặt đường.
    \item Khi mật độ tăng lên mức cao ($>25$ xe), \textbf{tất cả các kiến trúc đều chuyển sang trạng thái chệch âm (đếm thiếu phương tiện)} đối với xe máy do sự che khuất tương hỗ nghiêm trọng trong các dòng xe máy chen chúc.
    \item EfficientNet-B5 kiểm soát độ chệch âm tốt nhất (chỉ chệch $-0.63 \pm 0.87$ xe máy), trong khi ViT-Base/16 đếm thiếu trung bình $-4.63$ xe và ConvNeXt-Tiny đếm thiếu tới $-10.22$ xe. Đây chính là bằng chứng thực nghiệm giải thích nguồn gốc của sai số cảm biến mang tính hệ thống truyền vào Giai đoạn 2.
\end{itemize}

\subsection{Kiểm chứng trực quan hóa vùng chú ý bằng Grad-CAM++}
\label{sec:truc_quan_hoa_gradcam}
Để xác minh mô hình hồi quy đếm xe thực sự dựa trên các thuộc tính của phương tiện thay vì bám vào các tương quan giả tạo của cảnh nền, Hình \ref{fig:gradcam_chap5} minh họa bản đồ phân bổ đặc trưng Grad-CAM++ trích xuất từ tầng tích chập cuối của EfficientNet-B5.

\begin{figure}[htbp]
\centering
\includegraphics[width=0.92\linewidth]{fig/gradcam++.png}
\caption{Trực quan hóa bản đồ nhiệt Grad-CAM++ trên các khung hình camera giao thông, xác nhận vùng kích hoạt tập trung chuẩn xác vào thân xe ô tô và các luồng xe máy.}
\label{fig:gradcam_chap5}
\end{figure}

Bản đồ nhiệt Grad-CAM++ cho thấy các vùng kích hoạt cường độ cao (vùng màu đỏ và vàng) tập trung chính xác vào thân xe ô tô, mũ bảo hiểm và cụm đầu xe máy đang lưu thông trên làn đường. Các khu vực nền tĩnh như vỉa hè, biển báo, tòa nhà và mặt đường không có xe hoàn toàn không phát sinh kích hoạt. Điều này khẳng định mạng nơ-ron đã học được biểu diễn ngữ nghĩa thực chất của mật độ phương tiện.

% ==============================================================================
\section{Kết quả thực nghiệm Giai đoạn 2 (Graph Forecasting Benchmark)}
\label{sec:thuc_nghiem_giai_doan_2}

\subsection{Đối sánh đa chiều trên đồ thị 608 nút giao thông}
Bảng \ref{tab:stage2_benchmark_chap5} trình bày hiệu năng dự báo đa bước thời gian trung bình qua 5 seed ngẫu nhiên độc lập ($\mathcal{S} = \{42, 100, 2024, 22, 99\}$) trên toàn bộ 608 nút camera của mạng lưới đường bộ đô thị. Các mô hình được so sánh tại các mốc thời gian đại diện: $t+1$ (5 phút), $t+3$ (15 phút) và $t+6$ (30 phút phía trước).

\begin{table}[htbp]
\centering
\caption{Hiệu năng dự báo dòng giao thông Stage 2 qua 5 seed ngẫu nhiên ($\text{Mean} \pm \text{Std}$).}
\label{tab:stage2_benchmark_chap5}
\small
\resizebox{\textwidth}{!}{%
\begin{tabular}{lccccccc}
\toprule
\textbf{Kiến trúc mô hình} & \textbf{Số tham số} & \textbf{MAE Tổng} & \textbf{MAPE (\%)} & \textbf{RMSE} & \textbf{MAE ($t+1$)} & \textbf{MAE ($t+3$)} & \textbf{MAE ($t+6$)} \\
\midrule
STGCN Baseline \cite{stgcn} & 375K & $3.2703 \pm 0.0128$ & $11.44\% \pm 0.08\%$ & $4.4245 \pm 0.0194$ & $2.8430 \pm 0.0141$ & $3.2787 \pm 0.0117$ & $3.5350 \pm 0.0250$ \\
STGCN (Large) & 487K & $3.2805 \pm 0.0134$ & $11.53\% \pm 0.11\%$ & $4.4354 \pm 0.0156$ & $2.8498 \pm 0.0140$ & $3.2861 \pm 0.0162$ & $3.5497 \pm 0.0311$ \\
Graph WaveNet \cite{wavenet} & 549K & $3.2776 \pm 0.0049$ & $11.43\% \pm 0.03\%$ & $4.4486 \pm 0.0087$ & $\mathbf{2.8142 \pm 0.0088}$ & $3.2727 \pm 0.0048$ & $3.5800 \pm 0.0033$ \\
ASTGCN \cite{astgcn} & 1.533K & $3.4218 \pm 0.0115$ & $11.78\% \pm 0.08\%$ & $4.6509 \pm 0.0101$ & $2.9876 \pm 0.0262$ & $3.4296 \pm 0.0145$ & $3.6984 \pm 0.0104$ \\
STAEformer \cite{staeformer} & 568K & $3.2953 \pm 0.0184$ & $11.51\% \pm 0.12\%$ & $4.4555 \pm 0.0192$ & $2.8324 \pm 0.0237$ & $3.2924 \pm 0.0206$ & $3.5937 \pm 0.0275$ \\
MegaCRN \cite{megacrn} & 527K & $3.3433 \pm 0.0107$ & $11.55\% \pm 0.09\%$ & $4.5580 \pm 0.0182$ & $2.8436 \pm 0.0048$ & $3.3434 \pm 0.0111$ & $3.6660 \pm 0.0160$ \\
DSTAGNN \cite{dstagnn} & 911K & $3.3457 \pm 0.0169$ & $11.69\% \pm 0.08\%$ & $4.5499 \pm 0.0185$ & $2.8447 \pm 0.0162$ & $3.3373 \pm 0.0179$ & $3.6729 \pm 0.0189$ \\
iTransformer \cite{itransformer} & \textbf{150K} & $3.3720 \pm 0.0143$ & $11.65\% \pm 0.08\%$ & $4.5832 \pm 0.0216$ & $2.8613 \pm 0.0153$ & $3.3613 \pm 0.0163$ & $3.7205 \pm 0.0154$ \\
\midrule
\textbf{TA-STGCN (Đề xuất)} & \textbf{453K} & $\mathbf{3.2401 \pm 0.0108}$ & $\mathbf{11.32\% \pm 0.04\%}$ & $\mathbf{4.3920 \pm 0.0140}$ & $2.8346 \pm 0.0090$ & $\mathbf{3.2430 \pm 0.0106}$ & $\mathbf{3.4947 \pm 0.0153}$ \\
\bottomrule
\end{tabular}%
}
\end{table}

\begin{figure}[htbp]
\centering
\includegraphics[width=0.90\linewidth]{fig/val_mae_benchmark.pdf}
\caption{Đường cong hội tụ sai số Validation MAE qua các epoch huấn luyện trung bình qua 5 seed ngẫu nhiên. TA-STGCN thể hiện tiến trình hội tụ mượt mà và phương sai giữa các lần chạy rất thấp.}
\label{fig:val_mae_curves_chap5}
\end{figure}

\begin{figure}[htbp]
\centering
\includegraphics[width=0.92\linewidth]{fig/error_by_horizon.pdf}
\caption{Xu hướng tích lũy sai số (MAE và MAPE) theo tầm nhìn dự báo từ 5 phút ($t+1$) đến 30 phút ($t+6$). TA-STGCN duy trì tốc độ tích lũy sai số thoải nhất trên tầm nhìn xa.}
\label{fig:error_growth_chap5}
\end{figure}

\subsection{Đánh giá khách quan trên luồng dữ liệu sạch}
\begin{enumerate}
    \item \textbf{Độ chính xác tương đương trên dữ liệu lý tưởng:} Trên luồng dữ liệu không bị can thiệp nhiễu nhân tạo, TA-STGCN đạt sai số điểm tốt nhất ($\text{MAE} = 3.2401 \pm 0.0108, \text{MAPE} = 11.32\% \pm 0.04\%$), nhỉnh hơn Graph WaveNet ($3.2776$) và STGCN Baseline ($3.2703$) một khoảng nhỏ ($\approx 0.03 - 0.04$ MAE). Kích thước hiệu ứng Cohen's $d$ ở mức không đáng kể ($d = -0.024$ so với Graph WaveNet và $d = -0.032$ so với STGCN). Nghiên cứu khẳng định một cách trung thực rằng: \textbf{Dưới điều kiện dữ liệu sạch lý tưởng, các mô hình học sâu đồ thị hàng đầu đạt hiệu năng dự báo tương đương nhau.}
    \item \textbf{Hạn chế độ trôi dạt sai số ở tầm nhìn dài hạn ($t+6$):} Khi mở rộng tầm nhìn dự báo lên 30 phút ($t+6$), sự khác biệt kiến trúc trở nên rõ rệt hơn (Hình \ref{fig:error_growth_chap5}). Tại bước $t+6$, TA-STGCN đạt $\text{MAE} = 3.4947$, giảm được $0.0853$ MAE so với Graph WaveNet ($3.5800$) và $0.0403$ MAE so với STGCN ($3.5350$). Cơ chế chú ý thời gian đa đầu giúp kìm hãm sự tích lũy sai số trôi dạt dọc theo cửa sổ chuỗi thời gian.
    \item \textbf{Kiểm chứng dung lượng tham số:} Khi mở rộng dung lượng của STGCN gốc lên 487K tham số (phiên bản STGCN Large, cao hơn mức 453K của TA-STGCN), sai số dự báo không những không giảm mà tăng nhẹ lên $3.2805 \pm 0.0134$. Điều này chứng minh rằng ưu thế của TA-STGCN hoàn toàn bắt nguồn từ thiên kiến quy nạp (inductive bias) của cơ chế Self-Attention theo thời gian chứ không phải do việc gia tăng dung lượng tham số thuần túy.
\end{enumerate}

% ==============================================================================
\section{Phân tích lan truyền sai số theo mật độ ở Giai đoạn 2}
\label{sec:lan_truyen_sai_so_giai_doan_2}

Bảng \ref{tab:stage2_density_chap5} chi tiết hóa sai số dự báo tổng thể và độ chệch theo từng phân lớp phương tiện của cả 8 mô hình Stage 2 trên ba phân tầng mật độ giao thông thực tế.

\begin{table}[htbp]
\centering
\caption{Sai số dự báo và độ chệch phân tách ở Stage 2 theo các mức mật độ giao thông ($\text{Mean} \pm \text{Std}$).}
\label{tab:stage2_density_chap5}
\small
\resizebox{\textwidth}{!}{%
\begin{tabular}{llccc}
\toprule
\textbf{Kiến trúc mô hình} & \textbf{Chỉ số đánh giá} & \textbf{Mật độ Thấp ($<10$)} & \textbf{Mật độ TB ($10-25$)} & \textbf{Mật độ Cao ($>25$)} \\
\midrule
\multirow{3}{*}{\textbf{TA-STGCN (Đề xuất)}} & MAE Tổng thể & $0.9737 \pm 0.0066$ & $\mathbf{2.0024 \pm 0.0054}$ & $\mathbf{3.4837 \pm 0.0120}$ \\
 & Độ chệch Ô tô con & $-0.11 \pm 0.02$ & $+0.01 \pm 0.01$ & $\mathbf{-0.01 \pm 0.02}$ \\
 & Độ chệch Xe máy & $-0.07 \pm 0.01$ & $+0.01 \pm 0.02$ & $\mathbf{-0.03 \pm 0.06}$ \\
\midrule
\multirow{3}{*}{STGCN Baseline \cite{stgcn}} & MAE Tổng thể & $\mathbf{0.9696 \pm 0.0068}$ & $2.0126 \pm 0.0064$ & $3.5170 \pm 0.0136$ \\
 & Độ chệch Ô tô con & $-0.09 \pm 0.02$ & $+0.05 \pm 0.05$ & $+0.11 \pm 0.11$ \\
 & Độ chệch Xe máy & $-0.07 \pm 0.05$ & $+0.04 \pm 0.07$ & $+0.08 \pm 0.14$ \\
\midrule
\multirow{3}{*}{Graph WaveNet \cite{wavenet}} & MAE Tổng thể & $0.9997 \pm 0.0098$ & $2.0263 \pm 0.0100$ & $3.5212 \pm 0.0051$ \\
 & Độ chệch Ô tô con & $-0.08 \pm 0.02$ & $+0.03 \pm 0.03$ & $-0.02 \pm 0.07$ \\
 & Độ chệch Xe máy & $-0.07 \pm 0.03$ & $-0.04 \pm 0.06$ & $-0.01 \pm 0.12$ \\
\midrule
\multirow{3}{*}{ASTGCN \cite{astgcn}} & MAE Tổng thể & $1.0387 \pm 0.0237$ & $2.0758 \pm 0.0086$ & $3.6823 \pm 0.0133$ \\
 & Độ chệch Ô tô con & $-0.09 \pm 0.03$ & $-0.03 \pm 0.04$ & $-0.09 \pm 0.06$ \\
 & Độ chệch Xe máy & $-0.06 \pm 0.04$ & $-0.11 \pm 0.04$ & $-0.30 \pm 0.11$ \\
\midrule
\multirow{3}{*}{STAEformer \cite{staeformer}} & MAE Tổng thể & $0.9749 \pm 0.0113$ & $2.0092 \pm 0.0165$ & $3.5461 \pm 0.0207$ \\
 & Độ chệch Ô tô con & $-0.09 \pm 0.03$ & $+0.02 \pm 0.05$ & $-0.01 \pm 0.07$ \\
 & Độ chệch Xe máy & $-0.05 \pm 0.04$ & $+0.11 \pm 0.10$ & $+0.25 \pm 0.32$ \\
\midrule
\multirow{3}{*}{MegaCRN \cite{megacrn}} & MAE Tổng thể & $1.0219 \pm 0.0095$ & $2.0410 \pm 0.0023$ & $3.5976 \pm 0.0137$ \\
 & Độ chệch Ô tô con & $-0.09 \pm 0.02$ & $-0.05 \pm 0.02$ & $-0.08 \pm 0.04$ \\
 & Độ chệch Xe máy & $-0.10 \pm 0.04$ & $-0.12 \pm 0.11$ & $-0.21 \pm 0.29$ \\
\midrule
\multirow{3}{*}{DSTAGNN \cite{dstagnn}} & MAE Tổng thể & $1.0301 \pm 0.0118$ & $2.0448 \pm 0.0087$ & $3.5981 \pm 0.0186$ \\
 & Độ chệch Ô tô con & $-0.04 \pm 0.03$ & $+0.05 \pm 0.04$ & $+0.09 \pm 0.06$ \\
 & Độ chệch Xe máy & $-0.07 \pm 0.04$ & $+0.01 \pm 0.05$ & $+0.14 \pm 0.23$ \\
\midrule
\multirow{3}{*}{iTransformer \cite{itransformer}} & MAE Tổng thể & $1.0505 \pm 0.0104$ & $2.0448 \pm 0.0090$ & $3.6268 \pm 0.0152$ \\
 & Độ chệch Ô tô con & $-0.07 \pm 0.03$ & $+0.01 \pm 0.04$ & $-0.05 \pm 0.07$ \\
 & Độ chệch Xe máy & $-0.08 \pm 0.01$ & $-0.05 \pm 0.03$ & $-0.15 \pm 0.13$ \\
\bottomrule
\end{tabular}%
}
\end{table}

\begin{figure}[htbp]
\centering
\includegraphics[width=0.90\linewidth]{fig/stage2_density_mae.pdf}
\caption{Sai số dự báo MAE của các mô hình đồ thị Stage 2 theo các phân tầng mật độ phương tiện thực tế (Error Propagation).}
\label{fig:stage2_density_mae_chap5}
\end{figure}

\textbf{Quy luật lan truyền sai số nhận dạng:}
Tại điều kiện mật độ thấp, tất cả các mô hình đồ thị đều hoạt động tương đồng (MAE $\approx 0.97 - 1.05$). Tuy nhiên, tại mức mật độ cao ($>25$ xe/nút), nơi sai số nhận dạng đếm thiếu của Giai đoạn 1 lớn nhất, sự khác biệt trở nên rõ nét. TA-STGCN đạt sai số MAE thấp nhất ($3.4837 \pm 0.0120$), đồng thời kiểm soát độ chệch xe máy về mức gần như triệt tiêu ($-0.03 \pm 0.06$ xe), trong khi STAEformer chệch dương $+0.25$ xe và ASTGCN chệch âm $-0.30$ xe. Điều này chứng minh cơ chế Self-Attention thời gian trong TA-STGCN đã đóng vai trò làm mượt và bù trừ hiệu quả độ chệch âm từ thị giác máy tính truyền lên.

% ==============================================================================
\section{Phân tích phương sai và kiểm định ý nghĩa thống kê}
\label{sec:phan_tich_phuong_sai_thong_ke}

Để kiểm chứng nghiêm cẩn độ tin cậy của các kết quả thực nghiệm, Bảng \ref{tab:seed_variance} báo cáo kết quả kiểm định thống kê chuyên sâu: Khoảng tin cậy Bootstrap 95\% (với 1.000 lần lấy mẫu lại), Độ lệch chuẩn ($\sigma$), Hệ số biến thiên ($\text{CV}$), kiểm định cặp Wilcoxon Signed-Rank kèm hiệu chỉnh Holm-Bonferroni và quy mô ảnh hưởng Cohen's $d_z$ trên toàn bộ hơn $1.8 \times 10^7$ điểm dự báo.

\begin{table}[htbp]
\centering
\caption{Kiểm định thống kê Wilcoxon Signed-Rank, khoảng tin cậy Bootstrap 95\% và quy mô ảnh hưởng Cohen's $d_z$ cho Stage 2 (đối sánh với TA-STGCN).}
\label{tab:seed_variance}
\small
\resizebox{\textwidth}{!}{%
\begin{tabular}{lccccccc}
\toprule
\textbf{Kiến trúc mô hình} & \textbf{Mean MAE} & \textbf{95\% Bootstrap CI} & \textbf{Độ lệch chuẩn} & \textbf{CV (\%)} & \textbf{p-value} & \textbf{Cohen's $d_z$} & \textbf{Cải thiện (\%)} \\
\midrule
Graph WaveNet \cite{wavenet} & $3.2776$ & $3.2733 - 3.2819$ & $0.0049$ & $\mathbf{0.15\%}$ & $< 0.001$ & $-0.024$ & $+1.16\%$ \\
ASTGCN \cite{astgcn} & $3.4218$ & $3.4117 - 3.4319$ & $0.0115$ & $0.33\%$ & $< 0.001$ & $-0.109$ & $+5.61\%$ \\
STGCN Baseline \cite{stgcn} & $3.2703$ & $3.2591 - 3.2815$ & $0.0128$ & $0.39\%$ & $< 0.001$ & $-0.032$ & $+0.93\%$ \\
STAEformer \cite{staeformer} & $3.2953$ & $3.2792 - 3.3114$ & $0.0184$ & $0.56\%$ & $< 0.001$ & $-0.041$ & $+1.70\%$ \\
MegaCRN \cite{megacrn} & $3.3433$ & $3.3339 - 3.3527$ & $0.0107$ & $0.32\%$ & $< 0.001$ & $-0.071$ & $+3.19\%$ \\
DSTAGNN \cite{dstagnn} & $3.3457$ & $3.3309 - 3.3605$ & $0.0169$ & $0.51\%$ & $< 0.001$ & $-0.071$ & $+3.26\%$ \\
iTransformer \cite{itransformer} & $3.3720$ & $3.3595 - 3.3845$ & $0.0143$ & $0.42\%$ & $< 0.001$ & $-0.087$ & $+4.07\%$ \\
\midrule
\textbf{TA-STGCN (Đề xuất)} & $\mathbf{3.2401}$ & $\mathbf{3.2306 - 3.2496}$ & $\mathbf{0.0108}$ & $\mathbf{0.33\%}$ & --- & --- & --- \\
\bottomrule
\end{tabular}%
}
\end{table}

\textbf{Phân tích kiểm định thống kê:}
\begin{itemize}
    \item Kiểm định phi tham số Friedman omnibus trên toàn bộ tập mẫu đánh giá đạt giá trị thống kê $\chi^2 = 9.025,15$ ($p < 0.001$). Kiểm định Wilcoxon Signed-Rank hậu nghiệm đều cho $p < 0.001$. Tuy nhiên, với kích thước mẫu cực lớn ($> 1.8 \times 10^7$ điểm), giá trị $p$ luôn có xu hướng tiến về 0. Do đó, bằng chứng mang tính thực tế phải được xem xét qua quy mô ảnh hưởng Cohen's $d_z$.
    \item Trên dữ liệu sạch, các giá trị Cohen's $d_z$ đều nằm trong khoảng rất hẹp ($|d| < 0.11$), xác nhận rằng trên luồng đầu vào hoàn hảo, sự chênh lệch độ chính xác điểm giữa các mô hình đồ thị là không lớn.
    \item Hệ số biến thiên cực thấp của TA-STGCN ($\text{CV} = 0.33\%$) và khoảng tin cậy Bootstrap 95\% hẹp ($[3.2306 - 3.2496]$) chứng minh rằng hiệu năng của mô hình đề xuất \textbf{hoàn toàn ổn định, không phụ thuộc vào yếu tố ngẫu nhiên của seed khởi tạo}, duy trì tính nhất quán vượt trội hơn so với các mô hình Transformer nặng như STAEformer ($\text{CV} = 0.56\%$).
\end{itemize}

% ==============================================================================
\section{Nghiên cứu độ bền vững trước nhiễu nhận dạng (Robustness Analysis)}
\label{sec:nghien_cuu_do_ben_vung_nhieu}

Đây là nội dung thực nghiệm cốt lõi chứng minh tính ưu việt của kiến trúc đề xuất. Để mô phỏng chân thực các tác động của môi trường giao thông phức tạp (chói sáng, rung lắc camera, che khuất ngẫu nhiên), nghiên cứu tiêm nhiễu phân phối Gamma bất đối xứng vào chuỗi thời gian đầu vào tại 4 cấp độ ứng suất:
\begin{itemize}
    \item \textbf{Mức 1 ($\Delta\text{MAE} = 3.49$):} Mô phỏng điều kiện camera hoạt động thuận lợi ban ngày.
    \item \textbf{Mức 2 ($\Delta\text{MAE} = 4.64$):} Tương đương mức sai số nhận dạng trung bình thực tế của Giai đoạn 1.
    \item \textbf{Mức 3 ($\Delta\text{MAE} = 10.25$):} Tương đương tình trạng tắc nghẽn nặng nề và che khuất cục bộ.
    \item \textbf{Mức 4 ($\Delta\text{MAE} = 13.95$):} Thử nghiệm ứng suất cực đoan ($\text{MAE}_{\max} + 2\sigma$).
\end{itemize}

\begin{figure}[htbp]
\centering
\includegraphics[width=0.88\linewidth]{fig/noise_robustness_study.pdf}
\caption{Đường cong suy thoái sai số MAE của các mô hình đồ thị dưới 4 mức độ nhiễu nhận dạng mô phỏng. TA-STGCN thể hiện năng lực kiểm soát suy thoái tốt nhất.}
\label{fig:noise_robustness_chap5}
\end{figure}

Bảng \ref{tab:bang_suy_thoai_nhieu_chap5} và Bảng \ref{tab:cohens_d_chap5} trình bày chi tiết sai số MAE và quy mô ảnh hưởng Cohen's $d$ hai mẫu độc lập (pooled Cohen's $d$) giữa các kiến trúc so với TA-STGCN.

\begin{table}[htbp]
\centering
\caption{Độ suy thoái sai số MAE của các mô hình đồ thị dưới 4 mức độ nhiễu nhận dạng ($\text{Mean} \pm \text{Std}$).}
\label{tab:bang_suy_thoai_nhieu_chap5}
\small
\resizebox{\textwidth}{!}{%
\begin{tabular}{lcccc}
\toprule
\textbf{Kiến trúc mô hình} & \textbf{Mức 1 (3.49)} & \textbf{Mức 2 (4.64)} & \textbf{Mức 3 (10.25)} & \textbf{Mức 4 (13.95)} \\
\midrule
STGCN Baseline \cite{stgcn} & $5.3093 \pm 0.1161$ & $6.6614 \pm 0.1532$ & $12.0733 \pm 0.3396$ & $13.1724 \pm 0.3515$ \\
Graph WaveNet \cite{wavenet} & $5.7470 \pm 0.1056$ & $7.2378 \pm 0.1324$ & $13.2723 \pm 0.3303$ & $15.1046 \pm 0.4398$ \\
ASTGCN \cite{astgcn} & $5.7709 \pm 0.1521$ & $6.9956 \pm 0.2346$ & $11.8790 \pm 0.5428$ & $13.4685 \pm 0.7556$ \\
STAEformer \cite{staeformer} & $\mathbf{5.2238 \pm 0.2964}$ & $6.4260 \pm 0.5436$ & $10.5916 \pm 1.3943$ & $10.8454 \pm 1.2209$ \\
MegaCRN \cite{megacrn} & $5.9061 \pm 0.1227$ & $7.3144 \pm 0.1412$ & $13.1555 \pm 0.5314$ & $15.1106 \pm 0.8717$ \\
DSTAGNN \cite{dstagnn} & $6.2166 \pm 0.0514$ & $7.8276 \pm 0.2608$ & $12.5753 \pm 0.6823$ & $13.3705 \pm 0.8442$ \\
iTransformer \cite{itransformer} & $5.7070 \pm 0.0999$ & $6.9084 \pm 0.1786$ & $10.3254 \pm 0.7281$ & $10.8422 \pm 0.9367$ \\
\midrule
\textbf{TA-STGCN (Đề xuất)} & $5.2401 \pm 0.1000$ & $\mathbf{6.3120 \pm 0.1932}$ & $\mathbf{9.6381 \pm 0.8461}$ & $\mathbf{10.1532 \pm 1.1153}$ \\
\bottomrule
\end{tabular}%
}
\end{table}

\begin{table}[htbp]
\centering
\caption{Quy mô ảnh hưởng Cohen's $d$ so với TA-STGCN dưới nhiễu nhận dạng (dương thể hiện TA-STGCN vượt trội).}
\label{tab:cohens_d_chap5}
\small
\begin{tabular}{lcc}
\toprule
\textbf{Kiến trúc đối chứng} & \textbf{Mức 2 (Nhiễu thực tế: 4.64)} & \textbf{Mức 4 (Nhiễu cực đoan: 13.95)} \\
\midrule
STGCN Baseline \cite{stgcn} & $+2.004$ & $+3.651$ \\
Graph WaveNet \cite{wavenet} & $+5.590$ & $+5.841$ \\
ASTGCN \cite{astgcn} & $+3.181$ & $+3.480$ \\
STAEformer \cite{staeformer} & $+0.279$ & $+0.592$ \\
MegaCRN \cite{megacrn} & $+5.924$ & $+4.953$ \\
DSTAGNN \cite{dstagnn} & $+6.604$ & $+3.253$ \\
iTransformer \cite{itransformer} & $+3.206$ & $+0.669$ \\
\bottomrule
\end{tabular}
\end{table}

\textbf{Luận điểm khoa học cốt lõi từ thực nghiệm nhiễu:}
\begin{enumerate}
    \item \textbf{Tính bất đối xứng về chi phí và độ bền vững (Robustness-vs-Cost Asymmetry):} Việc tích hợp cơ chế Self-Attention theo thời gian vào STGCN chỉ làm tăng nhẹ mức sai số không đáng kể trên dữ liệu sạch, nhưng mang lại một \textbf{bước nhảy vọt về độ bền vững dưới nhiễu}. Tại Mức 2 (mức sai số thực tế của Stage 1), TA-STGCN vượt trội STGCN với Cohen's $d = +2.004$; tại Mức 4, sai số của TA-STGCN chỉ là $10.1532 \pm 1.1153$ so với $13.1724 \pm 0.3515$ của STGCN Baseline ($d = +3.651$).
    \item \textbf{Sự phân hóa rõ nét giữa các dòng kiến trúc:} Các mô hình thuần tích chập hoặc hồi quy không có cơ chế chú ý (Graph WaveNet, MegaCRN, STGCN) nhanh chóng bị suy thoái nặng nề khi gặp nhiễu (MAE vọt lên $13.17 - 15.11$, Cohen's $d$ từ $+3.65$ đến $+5.84$). Trong khi đó, các mô hình tích hợp cơ chế chú ý thời gian (TA-STGCN, STAEformer, iTransformer) giữ vững sai số ở mức $10.15 - 10.84$.
    \item \textbf{Tầm quan trọng của trật tự xử lý (Temporal-First so với Spatial-First):} ASTGCN mặc dù có cơ chế Attention nhưng lại tính toán chú ý không gian trước, chú ý thời gian sau. Khi một nút camera gặp nhiễu đột biến, tầng Spatial Attention bị đánh lừa và lập tức khuếch tán giá trị sai lệch sang toàn bộ đồ thị mạng lưới. Ngược lại, TA-STGCN tính toán Temporal Attention trước trên từng nút độc lập, hoạt động như một \textbf{bộ lọc thông thấp thích ứng} triệt tiêu xung nhiễu trước khi tín hiệu được truyền qua toán tử tích chập không gian ChebNet.
    \item \textbf{Hiệu quả vượt trội về tài nguyên phần cứng:} TA-STGCN đạt độ bền vững ngang ngửa hoặc nhỉnh hơn STAEformer nhưng tiết kiệm tới $25.2\%$ bộ nhớ kích hoạt GPU đỉnh ($8.606,9$ MB so với $11.509,1$ MB) và cắt giảm $70.4\%$ số tham số so với ASTGCN ($453\text{K}$ so với $1.533\text{K}$).
\end{enumerate}

% ==============================================================================
\section{Phân tích tính giải thích của cơ chế chú ý thời gian}
\label{sec:kha_nang_giai_thich_attention}

Hình \ref{fig:attention_heatmap_chap5} và Bảng \ref{tab:attention_ty_le_chap5} trình bày phân tích định lượng ma trận trọng số chú ý thời gian $\mathbf{A} \in \mathbb{R}^{24 \times 24}$ trích xuất từ tầng Temporal Attention của TA-STGCN qua 5 khung giờ đặc trưng trong ngày.

\begin{figure}[htbp]
\centering
\includegraphics[width=0.92\linewidth]{fig/attention_comparison_evening_vs_noon.pdf}
\caption{So sánh ma trận trọng số chú ý thời gian $\mathbf{A} \in \mathbb{R}^{24 \times 24}$ giữa (a) Giờ cao điểm tối (16:30--18:30) và (b) Giờ bình thường buổi trưa (11:30--13:30).}
\label{fig:attention_heatmap_chap5}
\end{figure}

\begin{table}[htbp]
\centering
\caption{Định lượng phân bổ trọng số chú ý thời gian qua các khung giờ giao thông trong ngày ($\text{Mean} \pm \text{Std}$).}
\label{tab:attention_ty_le_chap5}
\small
\resizebox{\textwidth}{!}{%
\begin{tabular}{lcccc}
\toprule
\textbf{Khung thời gian} & \textbf{Lưu lượng TB} & \textbf{Trọng số gần (20m)} & \textbf{Trọng số xa (60--120m)} & \textbf{Tỷ lệ (Gần / Xa)} \\
\midrule
Đêm khuya (00:00--05:00) & 5.0 & $0.0427 \pm 0.0038$ & $0.0417 \pm 0.0016$ & $1.03\times \pm 0.13\times$ \\
Cao điểm sáng (07:30--09:30) & 35.5 & $0.0461 \pm 0.0021$ & $0.0389 \pm 0.0013$ & $1.19\times \pm 0.09\times$ \\
Buổi trưa (11:30--13:30) & 36.0 & $0.0416 \pm 0.0013$ & $0.0419 \pm 0.0011$ & $0.99\times \pm 0.06\times$ \\
\textbf{Cao điểm tối (16:30--18:30)} & \textbf{39.5} & $\mathbf{0.0662 \pm 0.0067}$ & $\mathbf{0.0304 \pm 0.0032}$ & $\mathbf{2.22\times \pm 0.41\times}$ \\
Đêm muộn (21:00--23:00) & 17.5 & $0.0644 \pm 0.0062$ & $0.0314 \pm 0.0023$ & $2.08\times \pm 0.35\times$ \\
\bottomrule
\end{tabular}%
}
\end{table}

\textbf{Cơ chế tái phân bổ trọng số thích ứng theo trạng thái giao thông:}
\begin{itemize}
    \item Trong các giai đoạn giao thông ổn định (như Buổi trưa 11:30--13:30), mô hình phân bổ trọng số gần như đồng đều tuyệt đối giữa các bước thời gian gần và xa (tỷ lệ $0.99\times \pm 0.06\times$, với giá trị trọng số xấp xỉ mức đều $1/24 \approx 0.0417$).
    \item Khi bước vào \textbf{Giờ cao điểm tối (16:30--18:30)} --- thời điểm lưu lượng xe máy tăng vọt và có biến động động học cực lớn --- mô hình tự động co hẹp tiêu điểm chú ý: tăng mạnh trọng số cho khoảng thời gian 20 phút gần nhất lên $0.0662 \pm 0.0067$ (tăng $58.8\%$ so với mức phân bổ đều) và nén trọng số lịch sử xa (60--120 phút trước) xuống $0.0304 \pm 0.0032$. Tỷ lệ Gần/Xa đạt đỉnh \textbf{$2.22\times \pm 0.41\times$}. Hiện tượng tương tự diễn ra vào giờ Đêm muộn ($2.08\times \pm 0.35\times$).
    \item Hành vi này chứng minh rằng cơ chế Self-Attention đã tự động học được cách chuyển dịch trạng thái: ưu tiên nắm bắt các xu hướng động học tức thời trong giờ cao điểm biến động và cân bằng toàn dải lịch sử trong giờ bình ổn.
\end{itemize}

% ==============================================================================
\section{Nghiên cứu loại trừ thành phần (Ablation Study)}
\label{sec:ablation_study}

Để cô lập và đánh giá chính xác vai trò đóng góp của từng thành phần cấu trúc trong TA-STGCN, Bảng \ref{tab:ablation_full_chap5} trình bày kết quả thử nghiệm triệt tiêu trên 10 biến thể kiến trúc độc lập qua 5 seed ngẫu nhiên.

\begin{table}[htbp]
\centering
\caption{Kết quả Ablation Study đánh giá các biến thể kiến trúc của TA-STGCN qua 5 seed ngẫu nhiên ($\text{Mean} \pm \text{Std}$).}
\label{tab:ablation_full_chap5}
\small
\resizebox{\textwidth}{!}{%
\begin{tabular}{lcccccc}
\toprule
\textbf{Biến thể mô hình} & \textbf{Số tham số} & \textbf{RAM đỉnh (MB)} & \textbf{MAE Tổng} & \textbf{MAE Ô tô} & \textbf{MAE Xe máy} & \textbf{MAPE (\%)} \\
\midrule
\textbf{TA-STGCN (Mô hình chính, $L=2$)} & \textbf{453K} & \textbf{8.606,9} & $\mathbf{3.2401 \pm 0.0108}$ & $\mathbf{1.2087 \pm 0.0019}$ & $\mathbf{2.7286 \pm 0.0096}$ & $\mathbf{11.32\% \pm 0.04\%}$ \\
STGCN Baseline (Không Attention) \cite{stgcn} & 375K & 4.837,6 & $3.2703 \pm 0.0128$ & $1.2188 \pm 0.0063$ & $2.7467 \pm 0.0191$ & $11.44\% \pm 0.08\%$ \\
STGCN (Large, 487K params) & 487K & 5.793,9 & $3.2805 \pm 0.0134$ & $1.2157 \pm 0.0032$ & $2.7620 \pm 0.0194$ & $11.53\% \pm 0.11\%$ \\
TA-STGCN ($h=1$, Đơn đầu) & 453K & 8.158,1 & $3.2501 \pm 0.0075$ & $1.2088 \pm 0.0025$ & $2.7357 \pm 0.0075$ & $11.32\% \pm 0.03\%$ \\
TA-STGCN ($C=32$, Kênh hẹp) & 126K & 10.591,5 & $3.2717 \pm 0.0057$ & $1.2177 \pm 0.0023$ & $2.7542 \pm 0.0047$ & $11.46\% \pm 0.05\%$ \\
TA-STGCN ($K=1$, Không đồ thị không gian) & 427K & 8.033,6 & $3.2488 \pm 0.0105$ & $1.2104 \pm 0.0017$ & $2.7366 \pm 0.0119$ & $11.32\% \pm 0.04\%$ \\
TA-STGCN ($h=8$, 8 đầu chú ý) & 453K & 9.205,4 & $3.2494 \pm 0.0118$ & $1.2103 \pm 0.0012$ & $2.7367 \pm 0.0118$ & $11.36\% \pm 0.04\%$ \\
TA-STGCN (Độ sâu $L=3$ khối) & 647K & 10.536,0 & $3.2539 \pm 0.0098$ & $1.2113 \pm 0.0019$ & $2.7404 \pm 0.0096$ & $11.38\% \pm 0.05\%$ \\
TA-STGCN (Attention đặt ở đầu khối) & 489K & 8.741,3 & $3.2485 \pm 0.0087$ & $1.2060 \pm 0.0018$ & $2.7301 \pm 0.0077$ & $11.33\% \pm 0.04\%$ \\
TA-STGCN (Attention đặt ở cuối khối) & 453K & 8.961,9 & $3.2680 \pm 0.0099$ & $1.2097 \pm 0.0020$ & $2.7468 \pm 0.0079$ & $11.44\% \pm 0.09\%$ \\
\bottomrule
\end{tabular}%
}
\end{table}

\textbf{Bốn phát hiện kiến trúc quan trọng từ Ablation Study:}
\begin{enumerate}
    \item \textbf{Vị trí đặt tầng Temporal Attention:} Cấu hình đặt Attention ở giữa hai tầng tích chập mang lại hiệu năng cao nhất ($\text{MAE} = 3.2401$), vượt trội so với đặt ở cuối khối ($3.2680$) hoặc ở đầu khối ($3.2485$). Việc đặt Attention ngay trước toán tử ChebNet bảo đảm các xung nhiễu thời gian tại từng nút được lọc sạch trước khi diễn ra sự lan truyền thông tin sang các nút lân cận trên đồ thị.
    \item \textbf{Cấu hình số đầu chú ý ($h=4$):} Số đầu chú ý $h=4$ đem lại kết quả tối ưu nhất ($3.2401$). Giảm xuống đơn đầu ($h=1$) làm sai số tăng lên $3.2501$ do mất khả năng nắm bắt đa dạng các tần số dao động; trong khi tăng lên $h=8$ khiến sai số tăng nhẹ ($3.2494$) và dung lượng RAM đỉnh tăng thêm $598.5$ MB do chi phí ma trận kích hoạt.
    \item \textbf{Số lượng khối không -- thời gian ($L=2$):} Tăng độ sâu mạng lên $L=3$ khối làm số lượng tham số tăng lên 647K và bộ nhớ đỉnh tăng lên 10.536 MB, nhưng MAE lại tăng lên $3.2539$. Điều này phản ánh hiện tượng làm mịn quá mức (over-smoothing) phổ biến trong các mạng nơ-ron đồ thị sâu.
    \item \textbf{Đóng góp của cấu trúc topo đồ thị không gian ($K=3$):} Khi triệt tiêu toán tử ChebNet (đặt $K=1$, bỏ qua hoàn toàn liên kết cạnh giữa các nút giao), sai số tăng lên $3.2488$. Kết quả này khẳng định cấu trúc topo kết nối đường bộ thực tế từ OpenStreetMap đóng góp thực chất vào độ chính xác dự báo lưu lượng toàn mạng lưới.
\end{enumerate}
